\chapter{Diets}

\section*{Definition and Overview}
The term "diet" has two meanings: the foods a person usually eats, and a specific eating plan followed to lose weight or manage a condition. There is a vast range of diets promoted for weight loss, health, and wellbeing, from balanced approaches such as the Mediterranean diet to restrictive and fad diets. The evidence consistently shows that the most effective and sustainable approaches are those that include a variety of nutritious foods, fit the individual's lifestyle, and can be maintained long-term. This chapter explains the different types of diets, how to choose one, and what to watch out for.

\section*{Epidemiology}
Dieting is common, especially among women, with many Australians attempting weight-loss diets each year. Most weight-loss diets result in short-term loss followed by regain, and yo-yo dieting is common. Chronic disease is strongly linked to diet, and therapeutic diets are important in managing diabetes, heart disease, kidney disease, and other conditions. Fad diets remain popular despite weak evidence, and nutrition misinformation is widespread online.

\section*{Aetiology and Pathophysiology}
\begin{itemize}
    \item \textbf{Balanced diets:} such as the Mediterranean diet and DASH diet, which emphasise vegetables, fruit, wholegrains, and healthy fats
    \item \textbf{Energy-restricted diets:} reduce total kilojoules for weight loss
    \item \textbf{Low-carbohydrate diets:} restrict carbohydrate foods, used for weight loss and diabetes
    \item \textbf{Plant-based diets:} vegetarian and vegan patterns, which can be very healthy when well planned
    \item \textbf{Very restrictive diets:} such as very low-calorie or elimination diets, which carry nutritional risks
    \item \textbf{Fad diets:} trendy plans with limited evidence, often promising rapid results
    \item \textbf{Therapeutic diets:} prescribed for medical conditions, such as low-salt, low-FODMAP, or gluten-free diets
\end{itemize}

\section*{Clinical Features}
\begin{itemize}
    \item An effective diet is sustainable, nutritious, and suited to the individual
    \item Weight loss of 0.5--1 kg per week is a healthy, achievable rate
    \item Most weight loss occurs from eating fewer kilojoules than the body uses
    \item The best diet is one the person can maintain long-term
    \item Warning signs of a poor diet: extreme restriction, promises of rapid weight loss, elimination of whole food groups, and costly products
    \item Therapeutic diets are designed for specific conditions and should be guided by professionals
\end{itemize}

\section*{Differential Diagnosis}
\begin{itemize}
    \item Evidence-based versus fad diets
    \item Weight-loss diets versus diets for managing a medical condition
    \item Healthy plant-based diets versus poorly planned vegan or vegetarian diets
    \item Restrictive diets versus eating disorders, which need professional care
\end{itemize}

\section*{When to Seek Medical Care}
People considering a significant dietary change, especially for weight loss or a medical condition, should discuss it with a doctor or accredited dietitian. Anyone with a history of eating disorders, or who finds dieting is causing distress or obsessional thinking, should seek help.

\textbf{Call 000 (or local emergency services) for:}
\begin{itemize}
    \item Severe allergic reactions to food
    \item Choking or any food-related emergency
    \item Any medical emergency
\end{itemize}

\section*{Diagnosis and Investigations}
\begin{itemize}
    \item \textbf{Dietary assessment:} review of current eating patterns and goals
    \item \textbf{Health assessment:} weight, blood pressure, cholesterol, and blood glucose
    \item \textbf{Nutritional review:} blood tests for deficiencies where relevant
    \item \textbf{Evidence review:} assessment of the diet's safety and evidence base
    \item \textbf{Professional guidance:} referral to a dietitian for individualised plans
\end{itemize}

\section*{Management}
\begin{itemize}
    \item \textbf{Choose a balanced pattern:} base the diet on vegetables, fruit, wholegrains, lean protein, and healthy fats
    \item \textbf{Set realistic goals:} gradual weight loss is more sustainable
    \item \textbf{Make gradual changes:} small, lasting changes beat short-term overhaul
    \item \textbf{Watch portions and kilojoules:} portion control matters
    \item \textbf{Include enjoyable foods:} total restriction often backfires
    \item \textbf{Combine with activity:} exercise supports weight management and health
    \item \textbf{Seek professional advice:} for medical conditions and safe therapeutic diets
\end{itemize}

\section*{Complications}
\begin{itemize}
    \item Nutritional deficiencies from restrictive diets
    \item Weight regain and yo-yo dieting
    \item Disordered eating and poor relationship with food
    \item Muscle loss with crash dieting
    \item Worsening of medical conditions with inappropriate diets
    \item Financial cost of diet products and programs
\end{itemize}

\section*{Prognosis and Outcomes}
Evidence-based diets that are balanced and sustainable produce lasting health benefits, including weight management, better blood glucose and cholesterol, and reduced disease risk. Diets that work for one person may not suit another, so individualised guidance is important. The most successful "diet" is not a temporary plan but a long-term pattern of healthy eating that fits daily life.

\section*{Prevention and Self-Care}
\begin{itemize}
    \item Be wary of diets that promise rapid or extreme results
    \item Focus on adding healthy foods rather than only removing foods
    \item Eat mindfully and listen to hunger and fullness cues
    \item Cook at home more often and plan meals
    \item Seek advice from an accredited dietitian
    \item Avoid dieting in ways that cause distress or restriction guilt
    \item Combine healthy eating with activity and good sleep
\end{itemize}

\section*{Sources and Further Reading}
The following reputable sources provide further information for healthcare students and patients seeking reliable, current advice about diets.

\begin{itemize}
    \item Dietitians Australia. \textit{Choosing a diet: evidence-based advice.} https://dietitiansaustralia.org.au/
    \item Eat For Health (NHMRC). \textit{Australian Dietary Guidelines.} https://www.eatforhealth.gov.au/
    \item Healthdirect Australia. \textit{Weight loss and diets.} https://www.healthdirect.gov.au/
    \item NHS. \textit{Weight loss: safe and sustainable approaches.} https://www.nhs.uk/live-well/
    \item World Health Organization. \textit{Healthy diet fact sheet.} https://www.who.int/
    \item The Conversation. \textit{Examining popular diets: what the evidence says.} https://theconversation.com/
\end{itemize}
